\subsection*{Lisenssi}
%\includegraphics[scale=0.3]{commons/logo/cc.pdf} \includegraphics[scale=0.3]{commons/logo/by.pdf} \\

Tämän teoksen käyttöoikeutta koskee Creative Commons Nimeä 4.0 Kansainvälinen (CC BY 4.0) -lisenssi.
This work is licensed under a Creative Commons Attribution 4.0 International (CC BY 4.0) License.

\qrlinkki{http://creativecommons.org/licenses/by/4.0/deed.fi}{Lisenssin suomenkielinen tiivistelmä} \\
\qrlinkki{http://creativecommons.org/licenses/by/4.0}{Lisenssin englanninkielinen tiivistelmä} \\
\qrlinkki{http://creativecommons.org/licenses/by/4.0/legalcode}{Tarkat lisenssiehdot (englanniksi)}

Voit vapaasti:
\begin{description}
\item[Jakaa.] Kopioida ja jatkolevittää materiaalia, missä tahansa välineessä ja muodossa, mihin tahansa tarkoitukseen, myös kaupallisesti.
\item[Muuntaa.] Remiksata, muuntaa ja jalostaa materiaalia, mihin tahansa tarkoitukseen, myös kaupallisesti.
\end{description}

Lisenssinantaja ei voi vetää näitä vapauksia pois niin kauan kun seuraat lisenssin ehtoja.

Seuraavilla ehdoilla:
\begin{description}
\item[Nimeä.] Sinun on mainittava tekijät, näytettävä linkki lisenssiin sekä kerrottava, jos olet muuntanut materiaalia. Voit tehdä tämän millä tahansa kohtuullisella tavalla, mutta et niin, että lisenssin antaja näyttäisi tukevan sinua tai materiaalin käyttöäsi.
\item[No additional restrictions.] Et voi lisätä lakiehtoja tai teknisiä estoja, jotka estäisivät muita tekemästä mitään, minkä lisenssi sallii.
\end{description}

Huomautukset:
\begin{description}
\item[1.] Sinun ei tarvitse seurata lisenssiehtoja niiden materiaalin osien kohdalla, jotka ovat public domainissa tai joiden käyttö on sallittua tilannekohtaisen tekijänoikeuspoikkeuksen tai -rajoituksen vuoksi.
\item[2.] Mitään takuita ei ole annettu. Lisenssi ei välttämättä anna sinulle kaikkia käyttötarkoituksesi edellyttämiä oikeuksia. Esimerkiksi julkisuus-, yksityisyys- tai moraalisäädökset voivat rajoittaa materiaalin käyttöäsi.
\end{description}

\input{config/authors.tex}

